\subsection{Historical context and the universality prediction}

The exact solution of the nearest-neighbor Ising model made its critical
behavior a central example in statistical mechanics. Onsager determined
the square-lattice free energy and its singularity \cite{Onsager1944};
Yang obtained the spontaneous magnetization \cite{Yang1952}. The
conformal-field-theory framework of Belavin, Polyakov, and Zamolodchikov
\cite{BPZ1984} organizes the critical spin and energy fields by their
dimensions $1/8$ and $1$. The universality question asks whether the same
continuum correlations persist when a microscopic interaction destroys
the exact nearest-neighbor structure. In \Cref{model:main}, the
perturbation remains fixed as the lattice mesh tends to zero.

For the critical reference model, discrete complex analysis supplies a
rigorous route to conformal covariance. Chelkak and Smirnov developed
uniform discrete analytic estimates on isoradial graphs
\cite{CS2011} and proved convergence of critical Ising fermionic
observables in that setting \cite{CS2012}. Hongler and Smirnov computed
the one-point energy-density limit for fixed and free boundary
conditions \cite{HS2013}. Chelkak, Hongler, and Izyurov established
bulk spin-correlation convergence using discrete holomorphic spinors
\cite{CHI2015}, and subsequently treated mixed primary fields in
finitely connected domains with fixed, free, or mixed boundary arcs
\cite{CHI2022}. Dub\'edat's exact bosonization relates squared spin
correlations to neutral-charge expressions \cite{Dubedat2011}.
The bosonization theorem of Bayraktaro\u{g}lu, Izyurov, Virtanen, and Webb
\cite{BIVW2025} provides a modern formulation for primary fields in
multiply connected domains. These reference results identify the
continuum functions used here. Their exact hypotheses, and the
particular homogeneous spin formulas needed for the radial argument,
are recorded in \Cref{sec:inputs}.

Substantial universality results are also known beyond the
nearest-neighbor model. Giuliani, Greenblatt, and Mastropietro
\cite[Theorems~1.1--1.2]{GGM2012} proved the scaling limit of truncated energy correlations
for small, finite-range, rotation-invariant pair perturbations of either
sign. Their
construction first takes the thermodynamic limit at a noncritical
temperature depending on the mesh, then takes the scaling limit,
including the zero limiting-mass case. The interacting Grassmann
representation and constructive fermionic renormalization underlying
that work are important predecessors of rigorous perturbative Ising
universality, building on the Pinson--Spencer approach described in
\cite{GGM2012}.

Antinucci, Giuliani, and Greenblatt \cite[Theorem~1.1 and Corollary~1.2]{AGG2023} proved energy
universality on cylinders with periodic and free sides, with the
half-plane obtained by exhaustion. Their interaction already includes
finite-range even multispin terms of either sign; their energy is
centered by its actual interacting-domain expectation. They construct
analytic critical parameters and amplitudes and control the boundary
corrections, including the cancellation of a marginal boundary term.
Thus signed multispin interactions and actual-domain energy centering
are established features of that earlier energy theory. Cava,
Giuliani, and Greenblatt \cite[Theorem~1.1]{CGG2025} subsequently proved the
critical multipoint boundary-spin limit in the free half-plane, with
its multiplicative amplitude and the same Pfaffian limit as the
nearest-neighbor model.
Those boundary spins have dimension $1/2$; the bulk spins in
\Cref{model:main} have dimension $1/8$.

Giuliani's formulation of the universality program
\cite[Section~2 and Open Problems~1--2]{GiulianiICM2022} asks for
mixed spin and energy limits with two field normalizations under weak
even finite-range perturbations. \Cref{model:main} establishes this
prediction for the square-symmetric model in smooth simply connected
domains with free or either fixed boundary, and in the plane. It treats
all separated mixed bulk configurations with one temperature branch
and one pair of amplitudes. The broader program also includes other
topologies and boundary observables; those are outside the statement
proved here.

\subsection{Obstacles and main ideas}

Write $\E_0$ for expectation in the critical nearest-neighbor reference
law of the geometry under consideration. Our first task is a uniform
estimate for conditional averaging in the plane. Let $F$ be a bounded function
of the spins in a square of radius $s$. Conditioning on a square contour
at radius $R\gg s$ gives a function of the cutting spins. We show that,
after subtracting its reference mean, this function has essentially one
leading component in each spin-flip sector: the spin component for odd
$F$, and the energy component for even $F$. A distant fixed-plus disk
detects the coefficient of that component. The remainder decays faster
as $R/s$ grows. These estimates must hold uniformly over all bounded
$F$, including functions of arbitrarily many microscopic spins.

This uniformity is the first obstacle. Convergence of each fixed list
of primary fields gives information on finitely many prescribed spin
products. The local interaction functions created by a change of scale
need not belong to any such fixed list. The fermionic approach to energy
universality uses local representations of its observables, whereas
bulk spin insertions introduce nonlocal defects; see
\cite{GGM2012,AGG2023}. We instead use local functions of the physical
spins throughout.

The companion \emph{Buffered comparison and stopping-band resolution
in critical Ising} \cite{IsingToolkit} supplies comparison and stopping
estimates for the approximation step. Using these estimates, we show
that for every prescribed error, the effect of $F$ on all distant
bounded tests can be replaced by a finite signed combination of reference
laws with deterministic signed contours pinned. The total absolute
coefficient sum is at most $\|F\|_\infty$. The approximation follows by
conditioning on a separating contour and sorting a finite stopping
exploration into finitely many classes. When an outer continuation is
changed, its partition-function ratio provides the same bounded weight
for each inner alternative; this carries the comparison through the
unchanged collar.

We then replace the signed pins by normalized weights $e^X/\E_0e^X$,
where $X$ is a weighted sum of reference spins in thin tubes around the
contours. The pin-to-field comparison uses the companion's ordered
transmission estimate and continuity of reference correlations. The
reverse comparison adapts the cluster-necklace and ghost-attachment
argument of Camia, Jiang, and Newman \cite{CJN2020}. Magnetization moment
estimates \cite{CGN2015} justify these exponentials and successive
replacements at both ends of a pairing. These comparison arguments
all take place in the critical nearest-neighbor model.

Reflection positivity places the limiting soft weights in a radial
Hilbert space. Their norms are bounded independently of their profiles,
so the exact spin formulas give operator bounds for the whole finite
approximating family. The low spectrum contains the spin line of
dimension $1/8$ and the energy line of dimension $1$. In the even,
quarter-turn invariant sector, a finite charge identity cancels the
possible contribution of dimension two, leaving a remainder of order
$(s/R)^3$. This extra power is decisive: summing local interactions over
a two-dimensional block costs two powers of the scale ratio. The
remaining bulk directions therefore contract. Separate half-plane
estimates give the contraction needed for terms meeting the physical
boundary.

The second task is to turn conditional averaging into an exact change
of interactions. In a finite volume, represent an interaction as a list of bounded local
functions and expand its exponential in lists of occurrences. An
occurrence is conditionally averaged only when its support fits inside
its cutting square and an enlarged neighborhood of its support is
disjoint from those of all other occurrences. Overlapping occurrences are retained together. The
conditional Markov property preserves the reference integral of each
list. Taking the pointwise logarithm gives a new interaction $H'$ and
a scalar $c$. If $H$ is the sum of the input functions, then
\[
                 \E_0e^H=e^c\E_0e^{H'}.
\]
The connected logarithm and tree estimates give an analytic map with
exponential locality. We use the hard-core Mayer and tree-graph methods
\cite{FP2007,PY2017}, together with the occurrence-counting argument in
\Cref{sec:renormalization}. Exact local coordinate maps and contracted
remainders also occur in constructive renormalization
\cite{BS2015}; the reference kernels and list transformation used here
are proved directly.

The Kadanoff--Wilson scale-flow picture explains universality by the
decay of irrelevant interactions \cite{Kadanoff1966,Wilson1971}.
Here the linear estimate leaves one expanding coordinate, the thermal
one. Solving its recurrence backward while solving the contracting
coordinates forward expresses the thermal coordinate as an analytic
function of the contracting coordinates. This graph consists of
interactions that decay under repeated changes of scale. Related
sequence-space constructions appear in \cite{BBS2015}; our contraction argument is
\Cref{flow:stable-graph}. The microscopic temperature direction crosses
this graph transversely and selects $\beta_c(\lambda)$. Where a cutting
square lies in the interior of a physical domain, its conditional kernel
agrees with the plane kernel. Boundary contraction controls the remaining
edge terms; exponential locality controls the discrepancies on periodic
boxes. These comparisons carry the same temperature branch to all the
geometries in the theorem.

Finally we insert finitely many marked spin and energy observables in
the same exact partition identity. Each single source approaches a
conditional average of its reference field, multiplied by one
domain-independent amplitude. We
stop at a fixed small physical radius, while the source neighborhoods
are still disjoint. Connected interactions between distinct sources
then vanish, using exponential locality together with the decay of the
background or of a one-source remainder.

At this stopping scale the total residual background in a bounded domain
tends to zero, so the surviving source fields can be averaged in the
critical reference measure. This measure retains their mixed
correlations. Conditional independence in the disjoint cutting squares
recovers every subset moment of the microscopic reference fields.
Removing the energy singleton coefficients
from the logarithmic generating function gives centering by the actual
interacting state. An additional contraction for unscaled observables
constructs the periodic thermodynamic state before the continuum limit
and compares its residual local laws with the reference laws. The resulting massless correlations establish the critical behavior of
the selected branch.

\subsection{Proof structure}

\Cref{sec:inputs} fixes the reference conventions and states the exact
equilibrium interfaces and primary-field inputs.
\Cref{sec:buffers} proves the finite soft-field approximation, including
its uniformity over bounded tests and its moment bounds.
\Cref{sec:radial} constructs the radial space, proves the charge
cancellation and low-spectrum estimates, and derives the uniform
bulk and boundary conditional-averaging bounds.
\Cref{sec:renormalization} builds the exact analytic map and separates
the thermal, marked, and boundary directions.
\Cref{flow:section} constructs the temperature branch and compares its
plane trajectory with physical domains and periodic boxes.
\Cref{sec:correlations} tracks the marked sources, performs actual-state
energy centering, establishes the order of plane limits, and completes
\Cref{model:main}.

The order of choices is a large fixed scale ratio, a sufficiently
large initial lattice scale, and then a sufficiently small coupling
interval. No quantitative rate of primary-field or interface
convergence is required. Only the terminal physical scale depends on
the geometry of a fixed domain.
